\clearpage\fancyhead[L]{\small JUNO.}\SourceLeaf{397}
\begin{center}\textbf{GAUSS AN BODE.}\par\vspace{4pt}\hrule\vspace{4pt}\fontsize{10}{12}\selectfont Astronomisches Jahrbuch für 1819. Seite 219..221. Berlin 1816.\par\vspace{4pt}\hrule\vspace{4pt}\end{center}
\begin{flushright}Göttingen, den 10. August 1816.\end{flushright}

\par\indent Wie gewünscht, verehrtester Freund, sende ich Ihnen hiermit gerne die Ephemeriden für die nächste Erscheinung von Pallas und Juno und füge einige Anmerkungen hinzu.

\par\indent Die Ephemeride für Pallas hat Herr \textsc{Westphal} aus Schwerin berechnet, der sich hier zurzeit mit Eifer der Astronomie widmet. Dabei wurden die Störungen, wie seit 5 Jahren stets, nach meiner Theorie berücksichtigt. Daher ist die genaueste Übereinstimmung zu erwarten. Die nächste Opposition hat Herr \textsc{Westphal} wie folgt vorausberechnet:
\begin{center}\fontsize{10}{12}\selectfont
11. Juli 1817. $8^{\mathrm u}\;45^{\mathrm m}\;5^{\mathrm s}$ Mittlere Zeit in Göttingen\par
\setlength{\tabcolsep}{3pt}\begin{tabular}{lcr}
Wahre Länge & $.\;.\;.\;.\;.\;.\;.\;.\;.\;.\;.\;.\;.\;.$ & $289^\circ\;3'\;22''$\\
Geozentrische Breite & $.\;.\;.\;.\;.\;.\;.\;.\;.\;.\;.$ & $43\;33\;1$\\
\end{tabular}\end{center}

\par\indent Pallas wird im nächsten Jahr sehr lichtschwach sein und nur die 10. bis 11. Größenklasse haben. Hier ist eine Übersicht über ihre Lichtstärke bei den letzten Oppositionen und der kommenden:
\begin{center}\fontsize{10}{12}\selectfont\setlength{\tabcolsep}{11pt}\begin{tabular}{lr|lr}
1812 & $0.01666$ & 1816 & $0.05997$\\
1813 & $0.01475$ & 1817 & $0.01289$\\
1814 & $0.05476$ & 1818 & $0.02122$\\
\end{tabular}\end{center}

\Needspace{9\baselineskip}\par\indent Juno habe ich selbst zweimal mit dem Kreismikrometer beobachtet:
\begin{center}\fontsize{10}{12}\selectfont\setlength{\tabcolsep}{5pt}\begin{tabular}{l|c|c|c}
1816 & Mittlere Zeit & Rektaszension & \shortstack{Südliche\\Deklination}\\\hline
Juni 12 & $10^{\mathrm u}\;47^{\mathrm m}\;27^{\mathrm s}5$ & $251^\circ\;46'\;8''0$ & $3^\circ\;42'\;4''7$\\
13 & $10\;27\;34.7$ & $251\;33\;47.6$ & $3\;41\;11.7$\\
\end{tabular}\end{center}

\Needspace{11\baselineskip}\par\indent Bei dem ungewöhnlich schlechten Wetter im vergangenen Frühjahr konnten keine weiteren Beobachtungen gemacht werden. Daraus habe ich folgende Ergebnisse für die Opposition abgeleitet:
\begin{center}\fontsize{10}{12}\selectfont
3. Juni 1816. $4^{\mathrm u}\;20^{\mathrm m}\;20^{\mathrm s}$ Mittlere Zeit in Göttingen\par
\setlength{\tabcolsep}{3pt}\begin{tabular}{lcr}
Wahre Länge & $.\;.\;.\;.\;.\;.\;.\;.\;.\;.\;.\;.$ & $252^\circ\;51'\;39''2$\\
Geozentrische Breite & $.\;.\;.\;.\;.\;.\;.\;.\;.$ & $18\;34\;3.1$ Nördl.\\
\end{tabular}\end{center}

\Needspace{9\baselineskip}\par\indent Herr \textsc{Enke}, der diese Rechnung später wiederholte und dabei auch die Beobachtungen aus Wilna heranzog, fand
\begin{center}\fontsize{10}{12}\selectfont
3. Juni \quad $4^{\mathrm u}\;17^{\mathrm m}\;31^{\mathrm s}$ \quad Mittlere Zeit in Göttingen\par
$252^\circ\;51'\;32''6$\par
$18\;34\;14.6$
\end{center}

\Needspace{15\baselineskip}\par\indent Aus der Kombination der letzten vier Oppositionen leitete ich folgende Elemente ab:
\SourceLeaf{398}\begin{center}\fontsize{9}{11}\selectfont\setlength{\tabcolsep}{2pt}\renewcommand{\arraystretch}{1.15}\begin{tabular}{lcr}
Mittlere Länge 6. Sept. 1817. $0^{\mathrm u}$ in Göttingen & $.\;.$ & $8^\circ\;51'\;41''$\\
Perihel & $.\;.\;.\;.\;.\;.\;.\;.\;.\;.\;.\;.\;.\;.$ & $53\;16\;46$\\
Knoten & $.\;.\;.\;.\;.\;.\;.\;.\;.\;.\;.\;.\;.\;.\;.\;.$ & $171\;9\;41$\\
Neigung & $.\;.\;.\;.\;.\;.\;.\;.\;.\;.\;.\;.\;.\;.\;.$ & $13\;3\;49$\\
Exzentrizitätswinkel & $.\;.\;.\;.\;.\;.\;.\;.\;.\;.\;.$ & $14\;46\;0$\\
Tägliche mittlere tropische Bewegung & $.\;.\;.\;.$ & $812''40187$\\
Logarithmus der großen Halbachse & $.\;.\;.\;.\;.\;.$ & $0.4269223$\\
\end{tabular}\end{center}

\par\indent (Es ist wohl überflüssig zu erwähnen, dass diese Elemente allein für die Berechnung der nächsten Ephemeride berechnet wurden. Solange die Störungen noch nicht entwickelt sind, ist dieses Verfahren das bequemste und genaueste.)
\par\indent Aus diesen Elementen hat Herr Prof. \textsc{Harding} eine Ephemeride für die nächste Erscheinung berechnet, von der ich Ihnen hier eine Abschrift beilege. Dieselbe Rechnung hat auch Herr \textsc{Posselt} ausgeführt. Seine Ephemeride weicht von der hier vorliegenden überall nur in den Sekunden ab und wird in der Zeitschrift von Herrn \textsc{von Lindenau} abgedruckt. Nach der Rechnung von Herrn Professor \textsc{Harding} aus den oben angegebenen Elementen fällt die nächste Opposition auf den folgenden Zeitpunkt.
\begin{center}\fontsize{10}{12}\selectfont
6. Sept. 1817. $4^{\mathrm u}\;10^{\mathrm m}\;30^{\mathrm s}$ mittlere Zeit in Göttingen.\par
\setlength{\tabcolsep}{3pt}\begin{tabular}{lcr}
Wahre Länge & $.\;.\;.\;.\;.\;.\;.\;.\;.\;.\;.\;.\;.$ & $343^\circ\;37'\;36''$\\
Geozentrische Breite & $.\;.\;.\;.\;.\;.\;.\;.\;.\;.\;.$ & $3\;6\;9$ Nördl.\\
Die Lichtstärke 1816 & $.\;.\;.\;.\;.\;.\;.\;.\;.\;.\;.$ & $0.01747$\\
\hfill1817 & $.\;.\;.\;.\;.\;.\;.\;.\;.\;.\;.$ & $0.11398$\\
\end{tabular}\end{center}

\Needspace{18\baselineskip}\par\indent Erlauben Sie mir noch ein paar Bemerkungen über veränderliche Sterne. Dem von Herrn Prof. \textsc{Harding} entdeckten Stern in der Jungfrau schreibt Herr \textsc{Koch} im Jahrbuch für 1818 eine Periode von 295 Tagen zu, ohne anzugeben, worauf sie beruht. Diese Periode ist aber mit den Beobachtungen von Herrn Prof. \textsc{Harding} nicht vereinbar. — Einen zweiten veränderlichen Stern bemerkte Herr Prof. \textsc{Harding} im Jahr 1811 im Wassermann südlich unterhalb von $\omega^2$. Ausgehend von der sechsten Größenklasse verblasst der Stern allmählich, bis er nach ungefähr $2\frac12$ Monaten völlig unsichtbar ist. Seine Helligkeitsschwankungen scheinen überhaupt keiner regelmäßigen Periode zu folgen. Seine mittlere Position im Jahr 1800 war
\begin{center}\fontsize{10}{12}\selectfont\setlength{\tabcolsep}{3pt}\begin{tabular}{lcr}
Rektaszension & $.\;.\;.\;.\;.\;.\;.\;.\;.\;.\;.$ & $353^\circ\;21'\;48''5$\\
Jährliche. Ver. & $.\;.\;.\;.\;.\;.\;.\;.\;.\;.\;.\;.\;.\;.\;.$ & $+46.56$\\
Südliche Deklination & $.\;.\;.\;.\;.\;.\;.\;.\;.\;.$ & $16^\circ\;23'\;13''4$\\
Jährliche Abnahme & $.\;.\;.\;.\;.\;.\;.\;.\;.\;.\;.\;.\;.\;.$ & $19.88$.\\
\end{tabular}\end{center}
